%% 第 2 章：限制算子与商算子
%% 本文件由 tools/md2tex.py 自动生成，请勿直接编辑。

\chapter{限制算子与商算子}


\begin{jdefn}{限制算子}{r:3:1}
若 $U$ 是 $T$ 的不变子空间，则可以只研究 $T$ 在 $U$ 内部的作用：

\[
T|_U:U\to U.
\]

它称为 $T$ 在 $U$ 上的限制算子。

限制算子用于研究已经被分离出的部分。
\end{jdefn}

\begin{jdefn}{商算子}{r:3:2}
若 $U$ 是 $T$ 的不变子空间，则可在商空间 $V/U$ 上定义：

\[
\widetilde T:V/U\to V/U,
\]

\[
\widetilde T(v+U)=T(v)+U.
\]

它称为 $T$ 关于 $U$ 的商算子。
\end{jdefn}

\section{商算子为何良定义}

若：

\[
v+U=w+U,
\]

则：

\[
v-w\in U.
\]

因为 $U$ 对 $T$ 不变：

\[
T(v-w)\in U.
\]

于是：

\[
T(v)+U=T(w)+U.
\]

所以 $\widetilde T(v+U)$ 不依赖陪集代表元的选择，是良定义的。

\section{限制与商的意义}

\begin{jitem}
\item %
限制算子研究已经找到的不变部分；
\item %
商算子研究忽略该部分后，剩余自由度上的作用。
\end{jitem}

它们是上三角化、主分解与 Jordan 标准型归纳构造的重要工具。
